\section*{Chapitre \ref{ch:b2:reaction-enthalpy} --- Enthalpies de réaction}

\begin{solution}{exo:b2:reaction-enthalpy:1}
$\Delta_r H^\circ = 3(-393.51) + 4(-285.83) - (-104.7) - 5(0) =
\qty{-2219.2}{kJ/mol}$, négative~: la réaction est \omterm{def:b2:reaction-enthalpy:exothermic}{exothermique}. (L'enthalpie de
combustion mesurée est la même, \qty{-2219.2}{kJ/mol}.)
\end{solution}

\begin{solution}{exo:b2:reaction-enthalpy:2}
Méthane~: $890.3/16.04 = \qty{55.5}{kJ/g}$, soit \qty{55.5}{MJ} par
kilogramme. Dihydrogène~: $285.83/2.016 = \qty{141.8}{kJ/g}$,
\qty{141.8}{MJ/kg}~: environ 2.6 fois plus par kilogramme (mais bien moins par
litre, le gaz étant léger).
\end{solution}

\begin{solution}{exo:b2:reaction-enthalpy:3}
$\Delta\nu_{\text{gaz}} = 6 - 7.5 = -1.5$, donc $\Delta_r H = \Delta_r U +
\Delta\nu_{\text{gaz}}RT = -3264 - 1.5 \times 8.314 \times 298.15 \times
10^{-3} = \qty{-3267.7}{kJ/mol}$, en accord avec la valeur calculée à partir des
enthalpies de formation, \qty{-3267.6}{kJ/mol}.
\end{solution}

\begin{solution}{exo:b2:reaction-enthalpy:4}
La cible est la première réaction moins deux fois la seconde~:
$\Delta_r H^\circ = -393.5 - 2(-283.0) = \qty{+172.5}{kJ/mol}$. Elle est
\omterm{def:b2:reaction-enthalpy:exothermic}{endothermique}~: le coke chaud se refroidit en transformant le dioxyde de carbone en monoxyde
de carbone.
\end{solution}

\begin{solution}{exo:b2:reaction-enthalpy:5}
Liaisons rompues~: C--H et Cl--Cl, $416 + 243 = \qty{659}{kJ/mol}$~; liaisons
formées~: C--Cl et H--Cl, $350 + 432 = \qty{782}{kJ/mol}$~; estimation
\qty{-123}{kJ/mol}. À partir des enthalpies de formation~: $-83.7 - 92.3 + 74.9 =
\qty{-101.1}{kJ/mol}$. L'écart de \qty{22}{kJ/mol} vient de la valeur moyenne de C--Cl,
moyenne sur plusieurs molécules, et de la liaison C--H rompue dans le
méthane, plus forte que la moyenne des quatre.
\end{solution}

\begin{solution}{exo:b2:reaction-enthalpy:6}
À \qty{298}{K}~: $-241.83 + 285.83 = \qty{44.00}{kJ/mol}$. \omterm{thm:b2:reaction-enthalpy:kirchhoff}{Loi de Kirchhoff} avec
$\Delta C_p = 33.59 - 75.35 = \qty{-41.76}{J/(K.mol)}$~:
$44.00 - 0.04176 \times 101.85 = \qty{39.75}{kJ/mol}$ à \qty{400}{K}. Les
tables donnent $-242.85 + 282.59 = \qty{39.74}{kJ/mol}$~: les capacités
thermiques constantes sont excellentes sur cet intervalle.
\end{solution}

\begin{solution}{exo:b2:reaction-enthalpy:7}
Ionisation~: $4.341 \times 96.485 = \qty{418.8}{kJ/mol}$~; attachement~:
$-3.613 \times 96.485 = \qty{-348.6}{kJ/mol}$. \omterm{def:b2:reaction-enthalpy:lattice-enthalpy}{Enthalpie réticulaire} $= 436.7 +
89.0 + 121.3 + 418.8 - 348.6 = \qty{717}{kJ/mol}$, inférieure aux
\qty{787}{kJ/mol} du chlorure de sodium~: \ce{K+} est plus gros que \ce{Na+}, les
ions sont plus éloignés et s'attirent moins.
\end{solution}

\begin{solution}{exo:b2:reaction-enthalpy:8}
Capacité thermique du calorimètre~: $C = 2.00/0.418 = \qty{4.785}{kJ/K}$.
Chaleur libérée~: $4.785 \times 0.583 = \qty{2.789}{kJ}$, pour $\xi =
\qty{0.0500}{mol}$ (l'acide est limitant, la base en léger excès)~:
$\Delta_r H = -2.789/0.0500 = \qty{-55.8}{kJ/mol}$.
\end{solution}

\begin{solution}{exo:b2:reaction-enthalpy:9}
\ce{C8H18(l) + 25/2O2(g) -> 8CO2(g) + 9H2O(l)}~: $\Delta_r H^\circ =
8(-393.51) + 9(-285.83) + 250.3 = \qty{-5470.3}{kJ/mol}$~; avec $M =
\qty{114.2}{g/mol}$, \qty{47.9}{MJ/kg}. Dioxyde de carbone~: $8 \times 44.01 =
\qty{352.1}{g}$ pour \qty{5.470}{MJ}, soit \qty{64.4}{g/MJ}, contre
$44.01/0.8903 = \qty{49.4}{g/MJ}$ pour le méthane~: le méthane porte plus d'hydrogène par
carbone, et brûler l'hydrogène libère de la chaleur sans dioxyde de carbone.
\end{solution}

\begin{solution}{exo:b2:reaction-enthalpy:10}
Produits de \ce{CH4 + 2O2 -> CO2 + 2H2O(g)}~: $C_p = 37.13 + 2 \times 33.59 =
\qty{104.3}{J/K}$~; $q = \qty{802.3}{kJ}$~; $T_f = 298 + 802\,300/104.3 \approx
\qty{7990}{K}$, presque trois fois la valeur dans l'air, car aucun diazote ne
prend sa part de la chaleur. La flamme réelle dans le dioxygène est bien plus froide~: bien avant \qty{7990}{K}, le dioxyde de carbone et l'eau
se dissocient en monoxyde de carbone, dihydrogène, dioxygène et atomes, réactions
\omterm{def:b2:reaction-enthalpy:exothermic}{endothermiques} qui absorbent l'essentiel de la chaleur~; et les capacités thermiques croissent avec
$T$.
\end{solution}

\begin{solution}{exo:b2:reaction-enthalpy:11}
$\Delta_r H^\circ(700) = \Delta_r H^\circ(298) + \int_{298}^{700}(\Delta_r a +
\Delta_r b\,T)\,\dd T = -91.88 + [-60.5 \times 401.85 + 0.0270 \times
(700^2 - 298.15^2)] \times 10^{-3} = -91.88 - 24.31 + 10.83 =
\qty{-105.4}{kJ/mol}$, à \qty{0.2}{kJ/mol} près de la valeur des tables
(\qty{-105.2}{kJ/mol}), alors que l'estimation à $C_p$ constante s'en écartait de
\qty{4.5}{kJ/mol}.
\end{solution}

\begin{solution}{exo:b2:reaction-enthalpy:12}
$\Delta_r H^\circ = 2 \times 90.29 = \qty{+180.6}{kJ/mol}$ (\omterm{def:b2:reaction-enthalpy:exothermic}{endothermique}).
$\Delta_r C_p^\circ = 2(29.85) - 29.12 - 29.38 = \qty{1.2}{J/(K.mol)}$~: entre
\qty{298}{K} et \qty{2000}{K}, l'\omterm{def:b2:reaction-enthalpy:reaction-quantity}{enthalpie de réaction} varie donc
d'environ $1.2 \times 1702 = \qty{2}{kJ/mol}$, un pour cent. Une réaction
\omterm{def:b2:reaction-enthalpy:exothermic}{endothermique} est favorisée par une température élevée (\cref{ch:b2:equilibrium-shifts})~:
l'air froid ne produit presque pas de monoxyde d'azote, les flammes les plus chaudes en produisent
le plus~; c'est pourquoi la température de combustion est un levier contre ce
polluant.
\end{solution}

\begin{solution}{pb:b2:reaction-enthalpy:1}
\textbf{1.} \ce{C4H10(g) + 13/2O2(g) -> 4CO2(g) + 5H2O(l)}.
\textbf{2.} $\Delta_r H^\circ = 4(-393.51) + 5(-285.83) + 125.6 =
\qty{-2877.6}{kJ/mol}$.
\textbf{3.} Avec l'eau à l'état de vapeur, $4(-393.51) + 5(-241.83) + 125.6 =
\qty{-2657.6}{kJ/mol}$~; la différence, \qty{220.0}{kJ}, est l'enthalpie de
vaporisation de cinq moles d'eau (\qty{44.0}{kJ/mol} chacune).
\textbf{4.} $2877.6/58.12 = \qty{49.5}{kJ/g}$ et $2657.6/58.12 =
\qty{45.7}{kJ/g}$.
\textbf{5.} La seconde~: l'eau sort sous forme de vapeur et sa chaleur de
condensation est perdue.
\textbf{6.} $230 \times 45.7 = \qty{1.05e4}{kJ}$, environ \qty{10.5}{MJ}.
\textbf{7.} À volume constant, la chaleur reçue est $\Delta U$ (aucun
travail des forces de pression)~: la bombe mesure $\Delta_r U$.
\textbf{8.} $\Delta\nu_{\text{gaz}} = 4 - 1 - 6.5 = -3.5$ (l'eau est
liquide).
\textbf{9.} $\Delta_r U = \Delta_r H - \Delta\nu_{\text{gaz}}RT = -2877.6 +
3.5 \times 2.479 = \qty{-2868.9}{kJ/mol}$.
\textbf{10.} $\xi = 0.500/58.12 = \qty{8.60e-3}{mol}$~; chaleur $8.60 \times
10^{-3} \times 2868.9 = \qty{24.68}{kJ}$~; $\Delta T = 24.68/10.00 =
\qty{2.468}{K}$.
\textbf{11.} Chaleur $10.00 \times 2.461 = \qty{24.61}{kJ}$~: $\Delta_r U =
-24.61/(8.603 \times 10^{-3}) = \qty{-2860.6}{kJ/mol}$ et $\Delta_r H =
-2860.6 - 8.7 = \qty{-2869.3}{kJ/mol}$, \qty{0.29}{\percent} au-dessous des
tables en valeur absolue~: chaleur perdue par le calorimètre ou combustion
légèrement incomplète.
\textbf{12.} $n = 1000/18.02 = \qty{55.5}{mol}$~; $Q = 55.5 \times 75.35
\times 80 = \qty{334.5}{kJ}$.
\textbf{13.} Chaleur libérée $16.0 \times 45.72 = \qty{731.5}{kJ}$~; rendement
$334.5/731.5 = 0.46$.
\textbf{14.} Dans les gaz de combustion chauds qui s'échappent autour de la casserole, dans
l'air chauffé par la flamme et les parois de la casserole, et en rayonnement.
\textbf{15.} $230 \times 45.72 \times 0.457/334.5 = \qty{14.4}{L}$.
\textbf{16.} Les \qty{6.5}{mol} de \ce{O2} s'accompagnent de $6.5 \times 78.08/20.95 =
\qty{24.23}{mol}$ de \ce{N2} et de $6.5 \times 0.93/20.95 = \qty{0.289}{mol}$
de \ce{Ar}.
\textbf{17.} \qty{4}{mol} de \ce{CO2}, \qty{5}{mol} de \ce{H2O(g)},
\qty{24.23}{mol} de \ce{N2}, \qty{0.289}{mol} de \ce{Ar}.
\textbf{18.} La flamme est adiabatique et isobare, donc $\Delta H = 0$. Le long d'un
chemin formé de la réaction à \qty{298}{K} ($\Delta H_1 =
\qty{-2657.6}{kJ}$) suivie de l'échauffement de ces produits, $\Delta H_2 =
-\Delta H_1 = \qty{+2657.6}{kJ}$~: toute la chaleur passe dans les gaz.
\textbf{19.} $C_p = 4(37.13) + 5(33.59) + 24.23(29.12) + 0.289(20.79) =
\qty{1028}{J/K}$.
\textbf{20.} $T_f = 298 + 2\,657\,600/1028 = \qty{2883}{K}$.
\textbf{21.} $T_f = 2300 + 100 \times (2657.6 - 2515.7)/(2655.7 - 2515.7) =
\qty{2401}{K}$.
\textbf{22.} Les capacités thermiques des gaz croissent avec la température (d'environ
un quart entre \qty{298}{K} et \qty{2400}{K})~: les valeurs à température ambiante
chauffent les gaz à trop bon compte.
\textbf{23.} Au-dessus de \qty{2000}{K}, le dioxyde de carbone et l'eau se dissocient en partie
(processus \omterm{def:b2:reaction-enthalpy:exothermic}{endothermique}), la flamme rayonne, et le mélange avec l'air n'est jamais parfait.
\textbf{24.} La \omterm{def:b2:reaction-enthalpy:flame-temperature}{température de flamme adiabatique} du butane dans l'air est
$\boldsymbol{T_{\text{ad}} \approx \qty{2.40e3}{K}}$.
\end{solution}
